% from ega3-1-fr.tex
\oldpage[0\textsubscript{III}]{74}
$\mathrm B_r^{pq}(A)=\mathrm B_r^{pq}(\mathrm T(X^\bullet))$ et
$\mathrm Z_r^{pq}(A)=\mathrm Z_r^{pq}(\mathrm T(X^\bullet))$ sont
stationnaires, donc la suite spectrale précédente est \emph{birégulière}
(\hyperref[0.11.1.3-fr]{11.1.3}). Nous noterons cette suite
$\mathrm E(A)=(\mathrm E_r^{pq}(A))$ et nous dirons que c'est la
\emph{suite spectrale du foncteur $\mathrm T$ relative à l'objet filtré
$A$}.
\end{env}

\begin{env}[13.6.5]
\label{0.13.6.5-fr}
Supposons maintenant vérifiées les hypothèses de
(\hyperref[0.13.6.3-fr]{13.6.3}), dont nous conservons les notations.
Comme $\mathrm F^i(X^\bullet)$ (resp. $\mathrm F^i(Y^\bullet)$) est
facteur direct de $X^\bullet$ (resp. $Y^\bullet$), on a
\[
(\mathrm T v)(\mathrm T(\mathrm F^i(X^\bullet)))
\subset\mathrm T(\mathrm F^{i+s}(Y^\bullet))
\]
pour tout $i\in\mathbb Z$~; les définitions de
(\hyperref[0.11.2.2-fr]{11.2.2}) montrent alors que pour
$1\leq r\leq+\infty$, $\mathrm T v$ définit un morphisme
$\mathrm B_r^{pq}(\mathrm T(X^\bullet))\to
\mathrm B_r^{p+s,q-s}(\mathrm T(Y^\bullet))$ et un morphisme
$\mathrm Z_r^{pq}(\mathrm T(X^\bullet))\to
\mathrm Z_r^{p+s,q-s}(\mathrm T(Y^\bullet))$, d'où un morphisme
\[
w_r:\mathrm E_r^{pq}(A)\to\mathrm E_r^{p+s,q-s}(B)~;
\]
de même, on a pour l'aboutissement des morphismes
$u_n:\mathrm R^n\mathrm T(A)\to\mathrm R^n\mathrm T(B)$ tels que
$u_n(\mathrm F^p(\mathrm R^n\mathrm T(A)))\subset
\mathrm F^{p+s}(\mathrm R^n\mathrm T(B))$.

La définition des $d_r^{pq}$ (M, XV, 1) montre en outre que les
diagrammes
\[
\xymatrix{
\mathrm E_r^{pq}(A)\ar[r]^{d_r^{pq}}\ar[d]_{w_r}
&\mathrm E_r^{p+r,q-r+1}(A)\ar[d]^{w_r}\\
\mathrm E_r^{p+s,q-s}(B)\ar[r]_{d_r^{p+s,q-s}}
&\mathrm E_r^{p+r+s,q-r-s+1}(B)
}
\]
sont commutatifs~; on en déduit un diagramme commutatif analogue pour
les isomorphismes $\alpha_r^{pq}$, que nous laisserons au lecteur le
soin d'expliciter. Enfin (\emph{loc. cit.}), on a aussi des diagrammes
commutatifs pour les aboutissements
\[
\xymatrix{
\mathrm E_\infty^{pq}(A)\ar[r]^{\beta^{pq}}\ar[d]_{w_\infty}
&\operatorname{gr}^p(\mathrm R^{p+q}\mathrm T(A))
  \ar[d]^{u_{p+q}}\\
\mathrm E_\infty^{p+s,q-s}(B)\ar[r]_{\beta^{p+s,q-s}}
&\operatorname{gr}^{p+s}(\mathrm R^{p+q}\mathrm T(B)).
}
\]
\end{env}

\begin{env}[13.6.6]
\label{0.13.6.6-fr}
Supposons en particulier qu'il existe un anneau $S$, muni d'une
filtration $(\mathrm F^i(S))_{i\in\mathbb Z}$, et un homomorphisme
d'anneaux
\begin{equation}
\label{0.13.6.6.1-fr}
\tag{13.6.6.1}
h:S\to\operatorname{Hom}_C(A,A)
\end{equation}
